% Generated by `taskgen course examples` from dlp.ic_individual_log@signed (seed 2).
% Do not edit: change examples.yaml and regenerate.
\begin{example}[второй этап, логарифм $-1$]\label{ex:ic-phase2-signed}
Пусть $p = 281$, $g = 230$ --- первообразный корень по модулю $p$, факторная база $F = \{-1, 2, 3, 5\}$ (вычеты берутся абсолютно наименьшими, $|r| < p/2$), и первый этап метода исчисления индексов уже дал логарифмы $L_{2} = \log_g 2 = 132$, $L_{3} = \log_g 3 = 83$, $L_{5} = \log_g 5 = 38$. Найти $x = \log_g 248$, т.\,е. решить \[230^x \equiv 248\ (\textrm{mod}\ 281).\]

\par\noindent\textit{Решение.} Перебираем $s = 0, 1, 2, \dots$, пока $h g^{s} \bmod p$ не разложится над $F$ (каждое следующее значение --- предыдущее, умноженное на $g = 230$):\[\begin{array}{r|r|r|l}s & h g^s \bmod p & r & \text{разложение} \\ \hline 0 & 248 & -33 & -3 \cdot 11 \\ 1 & 278 & -3 & -3 \\ \end{array}\]Логарифм $-1$: $g^{(p-1)/2} \equiv -1$, так как $-1$ --- единственный элемент порядка $2$ в $\mathbb{Z}_p^*$, а $g^{(p-1)/2}$ имеет порядок $2$; значит, $L_{-1} = 140$. Итак, $h g^{1} \equiv -3$, откуда $x + 1 \equiv L_{-1} + L_{3}\pmod{280}$:\[x \equiv 140 + 83 - 1 = 222\ (\textrm{mod}\ 280).\]Проверка: $230^{222} \equiv 248\ (\textrm{mod}\ 281)$.\[x = 222.\]
\end{example}
